\section{Final Results or Outputs}
\begin{justify}
The completed system accepts task data, analyses estimation quality, stores records in MySQL, and returns a meaningful response to the user. For example, if the expected time is 2 hours and the actual time is 3 hours, the system returns Underestimation because the task took longer than expected.
\end{justify}

\begin{table}[H]
\centering
\begin{tabular}{|M|M|M|M|}
\hline
\textbf{Expected Time} & \textbf{Actual Time} & \textbf{Estimation Type} & \textbf{Meaning} \\
\hline
2 & 3 & Underestimation & User planned less time than required. \\
\hline
4 & 2 & Overestimation & User planned more time than required. \\
\hline
3 & 3 & Accurate & User estimate matched actual time. \\
\hline
\end{tabular}
\caption{Sample Output Cases}
\end{table}

\section{Limitations Observed}
\begin{justify}
The current implementation uses hardcoded database credentials and hardcoded user id for task insertion. Passwords are stored in plain text in the current academic version. The AI logic is rule-based and does not yet learn from large historical data. These limitations can be improved in future versions.
\end{justify}
